\begingroup
\setlength{\parindent}{0pt}
\setlength{\parskip}{0.40\baselineskip}
\setlength{\abovedisplayskip}{6.5pt plus 1.5pt minus 1.5pt}
\setlength{\belowdisplayskip}{6.5pt plus 1.5pt minus 1.5pt}
\setlength{\jot}{5pt}
\clubpenalty=10000
\widowpenalty=10000

\section*{How the Proof Works}

We start with a smooth, incompressible flow of finite energy and keep its
viscosity fixed and positive. The problem is to prove that the flow remains
smooth for all time. Its velocity \(u\) and pressure \(p\) satisfy the
three-dimensional Navier--Stokes equation:
\[
 \begin{tikzpicture}[
  baseline=0pt,
  every node/.style={inner sep=0pt,outer sep=0pt},
  term/.style={anchor=base west},
  annotation/.style={anchor=base,font=\footnotesize},
  brace/.style={decorate,decoration={brace,mirror,amplitude=2.5pt},
   line width=0.3pt,line cap=round}
 ]
  \node[term] (acceleration) at (0,0) {$\displaystyle\partial_tu$};
  \node[term] (transport) at
   ($(acceleration.base east)+(54pt,0)$) {$\displaystyle(u\!\cdot\nabla)u$};
  \node[term] (pressure) at
   ($(transport.base east)+(54pt,0)$) {$\displaystyle-\nabla p$};
  \node[term] (diffusion) at
   ($(pressure.base east)+(54pt,0)$) {$\displaystyle\nu\Delta u$};
  \node[anchor=base] at
   ($(acceleration.base east)!0.5!(transport.base west)$) {$+$};
  \node[anchor=base] at
   ($(transport.base east)!0.5!(pressure.base west)$) {$=$};
  \node[anchor=base] at
   ($(pressure.base east)!0.5!(diffusion.base west)$) {$+$};
  \node[term] at ($(diffusion.base east)+(1.5pt,0)$) {$,$};
  \node[term] at ($(diffusion.base east)+(38pt,0)$)
   {$\displaystyle\nabla\!\cdot u=0.$};
  \foreach \term/\label in {
   acceleration/acceleration,
   transport/nonlinear transport,
   pressure/pressure,
   diffusion/viscous diffusion}{
   \draw[brace] ($(\term.base west)+(0,-6pt)$)
    -- ($(\term.base east)+(0,-6pt)$);
   \node[annotation] at ($(\term.base)+(0,-18pt)$) {\label};
  }
 \end{tikzpicture}
\]
When strain stretches a vortex, incompressibility makes the material tube
narrow as it lengthens. If circulation is retained, the narrowing makes it
spin faster. Viscosity spreads rotation into the surrounding fluid, so the
tube can narrow while losing circulation. Its vorticity
\(\omega=\nabla\times u\) satisfies
\[
 \partial_t\omega+(u\cdot\nabla)\omega
 =(\omega\cdot\nabla)u+\nu\Delta\omega.
\]
The hope is simple: viscosity smooths the fluid faster than nonlinear
transport can sharpen it.

A smaller core gives viscosity less distance to act across, but the faster
flow can also sharpen it more quickly. Under the Navier--Stokes rescaling
with \(\lambda>1\),
\[
 \begin{aligned}
 u_\lambda(x,t)&=\lambda u(\lambda x,\lambda^2t),\\
 p_\lambda(x,t)&=\lambda^2p(\lambda x,\lambda^2t),
 \end{aligned}
\]
transport and diffusion scale together at fixed \(\nu\). Smaller scales give
viscosity no automatic advantage.

The critical Sobolev order is \(s=\tfrac12\):
\[
 \|u_\lambda(\cdot,t)\|_{\dot H^s}
 =\lambda^{s-1/2}\|u(\cdot,\lambda^2t)\|_{\dot H^s}.
\]
Concentration makes the norm grow above that order and shrink below it;
at the critical order it stays invariant. Energy lies below this threshold.
The flow can become faster and more concentrated while its total energy
falls:
\[
 E[u_\lambda](t)=\lambda^{-1}E[u](\lambda^2t),
 \qquad E=\tfrac12\|u\|_2^2.
\]
An energy bound can survive the concentration we need to exclude.

Both transport and diffusion act over successively shorter times. If the
flow could keep repeating this concentration, those times would permit a
Zeno cascade at a finite \(T_*\):
\[
 \tau_j=\lambda^{-2j}\tau_0,
 \qquad \sum_{j=0}^{\infty}\tau_j<\infty.
\]
Scaling allows that possibility; the regularity
problem is whether the fluid can actually carry it out.

For a finite-time singularity, velocity would have to blow up. That requires
acceleration to blow up, then jerk, and so on, and so on. These changes in
velocity are governed by the Navier--Stokes equation. Differentiating it
shows what each further temporal order involves:
\[
 \begin{aligned}
 \partial_tu={}&\nu\Delta u-(u\cdot\nabla)u-\nabla p,\\
 \partial_t^2u={}&\nu\Delta\partial_tu
 -(\partial_tu\cdot\nabla)u\\
 &-(u\cdot\nabla)\partial_tu-\nabla\partial_tp,\\
 \partial_t^3u={}&\nu\Delta\partial_t^2u
 -(\partial_t^2u\cdot\nabla)u\\
 &-2(\partial_tu\cdot\nabla)\partial_tu
 -(u\cdot\nabla)\partial_t^2u
 -\nabla\partial_t^2p.
 \end{aligned}
\]
\needspace{10\baselineskip}
Each time derivative splits the transport product and applies a Laplacian
to the preceding response. Substituting the first equation into the second
exposes \(\Delta^2u\); the next substitution exposes \(\Delta^3u\).
Higher temporal orders reach further into the spatial derivatives, with
the nonlinear terms carried along:
\[
 \begin{aligned}
 \partial_t^{n+1}u={}&\nu\Delta\partial_t^nu\\
 &-\sum_{a=0}^{n}\binom na
 (\partial_t^au\cdot\nabla)\partial_t^{n-a}u
 -\nabla\partial_t^np.
 \end{aligned}
\]

At the critical height, this relation exposes the spatial regularity
behind those temporal demands. Write
\(E_s=\tfrac12\|\Lambda^su\|_2^2\), where
\(\Lambda=(-\Delta)^{1/2}\), and let \(H=E_{1/2}\).
Differentiating \(E_s\) and integrating the viscous term by parts moves
one spatial derivative onto each factor:
\[
 E_s'=-2\nu E_{s+1}+\mathcal P_s,
\]
with the transport and pressure contribution
\[
 \mathcal P_s=-\operatorname{Re}\langle\Lambda^su,
 \Lambda^s((u\cdot\nabla)u+\nabla p)\rangle_{L^2}.
\]
Starting at \(s=\tfrac12\), each further time derivative exposes the next
spatial energy:
\[
 \begin{aligned}
 H'&=-2\nu E_{3/2}+\mathcal P_{1/2},\\
 H''&=4\nu^2E_{5/2}-2\nu\mathcal P_{3/2}
        +\partial_t\mathcal P_{1/2}.
 \end{aligned}
\]
Continuing through \(H''',H^{(4)},\ldots\) reaches
\(E_{7/2},E_{9/2},\ldots\), together with the differentiated nonlinear
contributions.

\needspace{7\baselineskip}
Those spatial orders are enough for smoothness. Any chosen number \(k\)
of continuous spatial derivatives needs only a finite Sobolev height:
\[
 m>k+\tfrac32
 \quad\Longrightarrow\quad H^m(\mathbb R^3)\hookrightarrow C^k(\mathbb R^3).
\]
Every choice of \(k\) is covered by a sufficiently high half-integer order.
With finite energy, their intersection is
\[
 L^2\cap\bigcap_{n\ge0}\dot H^{n+1/2}
 =\bigcap_{m\ge0}H^m
 =H^\infty\hookrightarrow C^\infty.
\]
The temporal requirements of a singularity have led us into the spatial
requirements of smoothness. If these orders remain controlled together
through \(T_*\), their limiting velocity is smooth enough to start the
flow again.

If a singularity were to form, we would expect the opposite: smoothness at
every earlier instant, but an unbounded fixed Sobolev norm as \(t\) approaches
\(T_*\). The temporal responses would become unbounded too, even while total
energy remained finite. Having all finite derivatives before \(T_*\) would
not give us the control needed to carry them through it.

That control must come from the initial velocity
\(u_0\in H^\infty_\sigma(\mathbb R^3)\) under the equation.
Incompressibility fixes pressure through
\(-\Delta p=\partial_i\partial_j(u_i u_j)\), with its decay normalization.
Writing \(V_n=\partial_t^nu\) and \(Q_n=\partial_t^np\), the velocity
and pressure responses obey
\[
 \begin{aligned}
 V_{n+1}&=\nu\Delta V_n
 -\sum_{a=0}^n\binom na(V_a\cdot\nabla)V_{n-a}-\nabla Q_n,\\
 Q_n&=R_iR_j\sum_{a=0}^n\binom na V_{a,i}V_{n-a,j}.
 \end{aligned}
\]
\needspace{5\baselineskip}
Here \(R_i\) are the Riesz transforms and repeated spatial indices are
summed. Starting with \(V_0=u_0\) at \(t=0\), these relations determine
each subsequent response. Spatial differentiation fills out each temporal
row. The resulting collection of velocity and pressure derivatives is the
mixed jet generated by the datum.

We can retain any finite number of temporal rows and spatial orders.
Choosing a temporal depth \(N\) and spatial height \(M\) gives a rectangle
\(\mathcal R_{N,M}[u_0;T_*]\). The datum-generated rectangle theorem,
Theorem~\ref{thm:datum-rectangles}, controls its adjacent temporal rows
throughout the approach to \(T_*\):
\[
 \sup_{0<T<T_*}\sum_{n=0}^{N}\left(
 \|V_n\|_{L^2(0,T;H^{M+1})}^2
 +\|V_{n+1}\|_{L^2(0,T;H^{M-1})}^2
 \right)<\infty.
\]
The bound can depend on the orders \(N,M\), but remains finite as
\(T\) approaches \(T_*\). Since \(V_{n+1}=\partial_tV_n\), the
adjacent rows control both a response and its change in time. The spaces
\(H^{M+1}\) and \(H^{M-1}\) pair across \(H^M\); their time-integrable
product gives a continuous endpoint trace by the Hilbert-triple theorem:
\[
 V_n\in C([0,T_*];H^M).
\]

\needspace{11\baselineskip}
A larger rectangle retains the derivatives already present. On each
overlap, both rectangles contain derivatives of the original velocity and
pressure, so their entries and endpoint traces agree. We can increase
temporal depth and spatial height independently without losing that
agreement. Their compatible intersection gives
\[
 \begin{aligned}
 \mathcal I[u_0]
 &:=\bigl(\mathcal R_{N,M}[u_0;T_*]\bigr)_{N,M\ge0},\\
 u&\in\bigcap_{r,m\ge0}H^r(0,T_*;H^m(\mathbb R^3)).
 \end{aligned}
\]
Every finite request for derivatives fits inside one of these rectangles.
In particular, the velocity reaches \(T_*\) in every spatial Sobolev
space, and its traces agree as one limiting velocity:
\[
 \begin{aligned}
 u_*&=\lim_{t\uparrow T_*}u(t)\quad\text{in every }H^m,\\
 u_*&\in H^\infty_\sigma\hookrightarrow C^\infty.
 \end{aligned}
\]
Smooth local existence starts the flow again from \(u_*\), and uniqueness
joins it to the original solution. It continues beyond its supposed
maximal time. A finite \(T_*\) is impossible, so \(T_*=\infty\).

The constructed intersection now supplies finite estimates at any chosen
orders. For any finite time \(T\), temporal order \(r\), spatial order
\(k\), and \(2\le p\le\infty\), choose
\(m>k+\tfrac32-\tfrac3p\). Sobolev embedding gives
\[
 \partial_t^ru\in C([0,T];H^m)
 \hookrightarrow L^\infty(0,T;W^{k,p}).
\]
In particular, taking \(r=0\), \(k=1\), \(p=\infty\), and \(m=3\)
controls the total strain over any finite time:
\[
 \int_0^T\|\nabla u(t)\|_\infty\,dt
 \le T\sup_{0\le t\le T}\|\nabla u(t)\|_\infty<\infty.
\]
For a locally circular material tube with axial strain \(\sigma\),
\(\dot r/r=-\sigma/2\) and \(\sigma\le\|\nabla u\|_\infty\). Hence
\[
 r(t)\ge r(0)\exp\!\left(
 -\frac12\int_0^t\|\nabla u(s)\|_\infty\,ds
 \right)>0.
\]
Only a finite amount of stretching can accumulate in a finite time.
The tube cannot be squeezed to zero radius there.

\endgroup
